%============================================================
% CHƯƠNG 2: KHẢO SÁT VÀ PHÂN TÍCH YÊU CẦU
%============================================================
\chapter{KHẢO SÁT VÀ PHÂN TÍCH YÊU CẦU}
\label{chap:survey}

Chương này đi sâu vào việc phân tích yêu cầu thực tế của bài toán thông
qua việc khảo sát các hệ thống quan trắc môi trường hiện có trên thị
trường. Trên cơ sở khảo sát, các chức năng cốt lõi cần có trong hệ
thống sẽ được xác định rõ ràng và biểu diễn chi tiết thông qua biểu đồ
use case cùng với phần đặc tả chức năng tương ứng.

%------------------------------------------------------------
\section{Khảo sát hiện trạng}
\label{sec:survey:status}
%------------------------------------------------------------

\subsection{Tổng quan về nhu cầu quan trắc môi trường trong không gian
kín}
\label{subsec:survey:need}

Chất lượng không khí trong các không gian kín như văn phòng, phòng học,
nhà xưởng và tòa nhà thương mại đang nhận được sự quan tâm ngày càng
tăng. Theo các nghiên cứu, con người dành trung bình hơn 90\% thời gian
trong nhà, trong khi chất lượng không khí trong nhà có thể kém hơn
ngoài trời từ 2 đến 5 lần~\cite{who2021air}. Nhiệt độ và độ ẩm là hai
thông số cơ bản nhưng có ảnh hưởng trực tiếp đến sức khỏe, năng suất
làm việc và sự phát triển của vi sinh vật trong không gian.

Nhu cầu đặt ra cho một hệ thống quan trắc hiện đại không chỉ dừng ở
việc \textbf{đo và hiển thị} giá trị hiện tại mà còn phải:

\begin{itemize}[leftmargin=2cm]
  \item Giám sát \textbf{liên tục, tự động} tại nhiều điểm phân tán
    trong không gian mà không cần can thiệp thủ công.
  \item \textbf{Dự báo xu hướng} trong tương lai gần để người vận hành
    chủ động điều chỉnh trước khi vượt ngưỡng nguy hiểm.
  \item Cung cấp \textbf{cái nhìn không gian} về sự phân bố nhiệt độ
    và độ ẩm trong toàn bộ khu vực, không chỉ tại từng điểm đơn lẻ.
  \item \textbf{Cảnh báo tự động} khi các chỉ số vượt ngưỡng an toàn
    đã cài đặt.
  \item Có thể \textbf{mở rộng} số lượng node cảm biến dễ dàng với chi
    phí thấp.
\end{itemize}

\subsection{Các giải pháp hiện có trên thị trường}
\label{subsec:survey:existing}

Hiện nay, thị trường cung cấp nhiều giải pháp quan trắc môi trường ở
các mức độ khác nhau. Bảng~\ref{tab:compare_solutions} tổng hợp và so
sánh các nhóm giải pháp tiêu biểu.

\begin{table}[H]
  \centering
  \caption{So sánh các giải pháp quan trắc môi trường hiện có}
  \label{tab:compare_solutions}
  \begin{tabularx}{\textwidth}{>{\bfseries}p{3cm} X X X}
    \toprule
    Tiêu chí
      & Trạm quan trắc nhà nước
      & Thiết bị thương mại (Xiaomi, SenseAir)
      & Giải pháp DIY / Nghiên cứu \\
    \midrule
    Chi phí
      & Rất cao (>100 triệu/trạm)
      & Trung bình (500k--5 triệu/node)
      & Thấp (<500k/node) \\
    Mật độ triển khai
      & Thưa thớt
      & Vừa phải
      & Linh hoạt, tùy biến \\
    Khả năng dự báo
      & Có (cloud-based)
      & Hạn chế hoặc không có
      & Tùy thiết kế \\
    Phân bố không gian
      & Không
      & Không
      & Tùy thiết kế \\
    Hệ sinh thái
      & Đóng
      & Đóng, phụ thuộc vendor
      & Mở, tùy biến \\
    Phù hợp nghiên cứu
      & Không
      & Hạn chế
      & Cao \\
    \bottomrule
  \end{tabularx}
\end{table}

Các sản phẩm thương mại như Xiaomi Mi Air Monitor hay các cảm biến
SenseAir tuy có giao diện đẹp và dễ sử dụng, nhưng hoạt động trong hệ
sinh thái đóng, không cho phép tích hợp pipeline ML tùy chỉnh hay trực
quan hóa phân bố không gian. Các trạm quan trắc nhà nước có độ chính
xác cao nhưng chi phí vận hành rất lớn và mật độ quá thưa để phục vụ
giám sát trong không gian kín.

\subsection{Xu hướng công nghệ AIoT trong quan trắc môi trường}
\label{subsec:survey:trend}

Sự phát triển của hệ sinh thái IoT với các vi điều khiển như ESP32 có
tích hợp Wi-Fi và chi phí thấp đã mở ra hướng tiếp cận mới cho bài toán
quan trắc phân tán~\cite{shi2016edge}. Kết hợp với giao thức MQTT nhẹ
và hiệu quả~\cite{banks2014mqtt}, các node cảm biến có thể truyền dữ
liệu liên tục về hệ thống trung tâm với độ trễ thấp và tiêu thụ điện
năng tối thiểu.

Về phía xử lý dữ liệu, các mô hình học sâu cho chuỗi thời gian -
đặc biệt là LSTM (Long Short-Term Memory)~\cite{hochreiter1997lstm} -
đã chứng minh hiệu quả vượt trội so với các phương pháp thống kê
truyền thống như ARIMA trong bài toán dự báo dữ liệu cảm biến môi
trường. Kết hợp với các thuật toán nội suy không gian như
Kriging~\cite{cressie1990kriging} và IDW~\cite{shepard1968idw}, hệ
thống có thể tái tạo bản đồ phân bố môi trường liên tục từ một tập hợp
hữu hạn các điểm đo.

\subsection{Kết luận và đề xuất giải pháp}
\label{subsec:survey:conclusion}

Từ phân tích trên, nhận thấy khoảng trống rõ ràng trên thị trường: chưa
có giải pháp chi phí thấp, mã nguồn mở nào kết hợp được cả ba yếu tố:
(1) giám sát thời gian thực đa node, (2) dự báo chuỗi thời gian bằng
LSTM, và (3) trực quan hóa phân bố không gian bằng heatmap nội suy.

Đề tài đề xuất xây dựng hệ thống sử dụng vi điều khiển \textbf{ESP32}
kết hợp cảm biến \textbf{DHT11}, giao tiếp qua \textbf{MQTT}, backend
\textbf{Node.js/Express}, mô hình \textbf{LSTM} (PyTorch) phục vụ qua
\textbf{FastAPI}, và nội suy không gian \textbf{Kriging/IDW} để tạo
heatmap hiển thị trên \textbf{Leaflet}, nhằm lấp đầy khoảng trống đó.

%------------------------------------------------------------
\section{Tổng quan chức năng}
\label{sec:survey:usecase}
%------------------------------------------------------------

\subsection{Biểu đồ use case tổng quát}
\label{subsec:survey:uc:overall}

Hình~\ref{fig:usecase_overall} minh họa biểu đồ use case tổng quan của
hệ thống, thể hiện mối quan hệ giữa các tác nhân và các chức năng
chính. Hệ thống có hai tác nhân chính:

\begin{itemize}[leftmargin=2cm]
  \item \textbf{Người dùng (User):} Người vận hành hoặc quản lý không
    gian được giám sát. Tương tác chủ yếu qua giao diện web dashboard
    để xem dữ liệu, nhận cảnh báo và xem dự báo.
  \item \textbf{Node IoT (ESP32):} Thiết bị phần cứng đóng vai trò tác
    nhân tự động, định kỳ thu thập và gửi dữ liệu cảm biến lên hệ
    thống qua MQTT.
\end{itemize}

\begin{figure}[H]
  \centering
  \includegraphics[width=0.9\textwidth]{uc_overall}
  \caption{Biểu đồ use case tổng quan của hệ thống}
  \label{fig:usecase_overall}
\end{figure}

Các chức năng chính của hệ thống bao gồm: đăng ký và cấu hình node IoT,
thu thập dữ liệu cảm biến theo thời gian thực, xem dữ liệu và biểu đồ
lịch sử, xem dự báo chuỗi thời gian, xem heatmap phân bố không gian,
và nhận cảnh báo khi vượt ngưỡng.

\subsection{Biểu đồ use case phân rã: Đăng ký và cấu hình node IoT}
\label{subsec:survey:uc:register}

Hình~\ref{fig:usecase_register} mô tả quy trình khởi tạo và quản lý
node IoT trong hệ thống. Chức năng trọng tâm là quá trình
\textit{provisioning} tự động: node ESP32 khi khởi động lần đầu sẽ
publish một thông điệp lên topic \texttt{devices/register} kèm địa chỉ
MAC. Backend Aedes broker lắng nghe topic này, kiểm tra xem thiết bị đã
tồn tại chưa, nếu chưa thì tạo bản ghi trong cơ sở dữ liệu, sinh
\texttt{api\_key} và gửi cấu hình (ngưỡng cảnh báo, phòng được gán) về
node qua topic \texttt{devices/\{mac\_address\}/config}.

\begin{figure}[H]
  \centering
  \includegraphics[width=0.9\textwidth]{uc_register}
  \caption{Biểu đồ use case phân rã: Đăng ký và cấu hình node IoT}
  \label{fig:usecase_register}
\end{figure}

Người dùng Admin có thể quản lý thiết bị qua dashboard: xem danh sách
node, chỉnh sửa ngưỡng cảnh báo nhiệt độ, độ ẩm (cao/thấp), gán node
vào phòng/khu vực, và vô hiệu hóa node khi cần.

\subsection{Biểu đồ use case phân rã: Thu thập dữ liệu và dự báo}
\label{subsec:survey:uc:collect}

Hình~\ref{fig:usecase_collect} minh họa quy trình thu thập dữ liệu tự
động từ node IoT và pipeline xử lý phía backend.

\begin{figure}[H]
  \centering
  \includegraphics[width=0.9\textwidth]{uc_collect}
  \caption{Biểu đồ use case phân rã: Thu thập dữ liệu và dự báo}
  \label{fig:usecase_collect}
\end{figure}

Mỗi 5 giây, node ESP32 đọc dữ liệu từ cảm biến DHT11 và publish JSON
lên topic \texttt{device/\{api\_key\}/data}. Backend xác thực
\texttt{api\_key}, lưu vào bảng \texttt{readings} trong PostgreSQL,
đồng thời emit sự kiện \texttt{new\_reading} qua Socket.io đến các
client đang subscribe. Khi người dùng yêu cầu xem dự báo, backend gọi
FastAPI service, service này đọc 30 bản ghi gần nhất từ PostgreSQL,
đưa qua mô hình LSTM và trả về 15 giá trị dự báo tiếp theo.

\subsection{Biểu đồ use case phân rã: Xem dữ liệu và nhận cảnh báo}
\label{subsec:survey:uc:view}

Hình~\ref{fig:usecase_view} mô tả các chức năng tương tác chính của
người dùng với dashboard.

\begin{figure}[H]
  \centering
  \includegraphics[width=0.9\textwidth]{uc_view}
  \caption{Biểu đồ use case phân rã: Xem dữ liệu và nhận cảnh báo}
  \label{fig:usecase_view}
\end{figure}

Người dùng có thể xem dữ liệu realtime qua biểu đồ Chart.js được cập
nhật tức thì thông qua Socket.io, xem lịch sử dữ liệu theo khoảng thời
gian tùy chọn, xem kết quả dự báo 15 phút tiếp theo, và xem heatmap
phân bố nhiệt độ, độ ẩm trong không gian phòng trên bản đồ Leaflet. Khi
các chỉ số vượt ngưỡng đã cài đặt, hệ thống tự động ghi alert vào cơ
sở dữ liệu, hiển thị thông báo trên dashboard và gửi email cảnh báo.

\subsection{Quy trình nghiệp vụ}
\label{subsec:survey:workflow}

\subsubsection*{a. Quy trình thu thập và xử lý dữ liệu}

Hình~\ref{fig:workflow_data} mô tả quy trình nghiệp vụ chính khi node
IoT gửi dữ liệu về hệ thống. Quy trình diễn ra hoàn toàn tự động và
liên tục theo chu kỳ 5 giây:

\begin{enumerate}[leftmargin=2cm]
  \item Node ESP32 đọc nhiệt độ, độ ẩm từ DHT11 và sinh giá trị AQI mô phỏng.
  \item ESP32 đóng gói JSON gồm nhiệt độ, độ ẩm, AQI và publish lên MQTT broker
    (Aedes) qua topic \texttt{device/\{api\_key\}/data}.
  \item Backend lắng nghe sự kiện \texttt{publish} của Aedes, xác thực
    \texttt{api\_key} với PostgreSQL.
  \item Nếu hợp lệ, bản ghi được lưu vào bảng \texttt{readings}.
  \item \texttt{alertService} kiểm tra giá trị vừa nhận với ngưỡng
    cài đặt của thiết bị. Nếu vượt ngưỡng, tạo bản ghi trong bảng
    \texttt{alerts} và gửi email thông báo.
  \item Socket.io emit sự kiện \texttt{new\_reading} đến tất cả client
    đang lắng nghe room \texttt{device\_\{id\}} và
    \texttt{user\_\{id\}}.
  \item Dashboard React nhận sự kiện, cập nhật biểu đồ Chart.js mà
    không cần tải lại trang.
\end{enumerate}

\begin{figure}[H]
  \centering
  \includegraphics[width=\textwidth]{workflow_data}
  \caption{Quy trình nghiệp vụ thu thập và xử lý dữ liệu}
  \label{fig:workflow_data}
\end{figure}

\subsubsection*{b. Quy trình dự báo và nội suy không gian}

Hình~\ref{fig:workflow_forecast} mô tả quy trình khi người dùng yêu cầu
xem dự báo và heatmap phân bố không gian:

\begin{enumerate}[leftmargin=2cm]
  \item Người dùng mở trang dự báo hoặc trang heatmap trên dashboard.
  \item Frontend gọi REST API \texttt{GET /api/forecast} hoặc
    \texttt{GET /api/spatial-forecast} đến backend Node.js.
  \item Backend chuyển tiếp request tới FastAPI service: \texttt{forecast/main.py}
    cho dự báo trạm hoặc \texttt{forecast/spatial\_forecast.py} cho heatmap.
  \item FastAPI service truy vấn dữ liệu lịch sử 30 phút gần nhất từ PostgreSQL.
  \item Dữ liệu được chuẩn hóa qua \texttt{scaler.pkl}, đưa vào mô
    hình LSTM (PyTorch) để dự báo 15 phút tiếp theo.
  \item Với yêu cầu heatmap: kết hợp giá trị thực đo (hoặc dự báo) với
    tọa độ vị trí của từng node, áp dụng thuật toán IDW và Kriging để
    nội suy ra lưới dữ liệu 40$\times$50 điểm.
  \item Kết quả trả về dưới dạng JSON, frontend render lên bản đồ
    Leaflet dưới dạng heatmap màu sắc.
\end{enumerate}

\begin{figure}[H]
  \centering
  \includegraphics[width=\textwidth]{workflow_forecast}
  \caption{Quy trình nghiệp vụ dự báo và nội suy không gian}
  \label{fig:workflow_forecast}
\end{figure}

%------------------------------------------------------------
\section{Đặc tả chức năng}
\label{sec:survey:spec}
%------------------------------------------------------------

\subsection{Đặc tả use case UC01 - Đăng ký node IoT vào hệ thống}
\label{subsec:survey:uc01}

\begin{table}[H]
  \centering
  \caption{Đặc tả use case UC01: Đăng ký node IoT vào hệ thống}
  \label{tab:uc01}
  \begin{tabularx}{\textwidth}{p{4cm} X}
    \toprule
    \textbf{Mã use case}    & UC01 \\
    \midrule
    \textbf{Tên}            & Đăng ký node IoT vào hệ thống \\
    \textbf{Tác nhân}       & Node IoT (ESP32) \\
    \textbf{Mục đích}       & Cấp phát định danh và cấu hình ban đầu
                              cho node cảm biến mới \\
    \textbf{Tiền điều kiện} &
      \begin{minipage}[t]{\linewidth}
        \begin{enumerate}[leftmargin=*, topsep=0pt]
          \item ESP32 đã được cấp thông tin Wi-Fi qua AP cấu hình \texttt{Sensor-Setup-[MAC]}
          \item MQTT broker (Aedes) đang chạy và cho phép client provisioning kết nối
        \end{enumerate}
      \end{minipage} \\
    \textbf{Luồng chính}    &
      \begin{minipage}[t]{\linewidth}
        \begin{enumerate}[leftmargin=*, topsep=0pt]
          \item ESP32 khởi động, publish MAC address lên topic
            \texttt{devices/register}
          \item Aedes broker nhận message, kiểm tra MAC trong DB
          \item Nếu chưa tồn tại: tạo bản ghi device, sinh
            \texttt{api\_key} ngẫu nhiên, lưu vào bảng
            \texttt{devices}
          \item Backend publish cấu hình (ngưỡng, \texttt{api\_key})
            về topic \texttt{devices/\{mac\}/config}
          \item ESP32 nhận cấu hình, lưu vào Flash qua
            \texttt{Preferences.h}
        \end{enumerate}
      \end{minipage} \\
    \textbf{Luồng thay thế} &
      \begin{minipage}[t]{\linewidth}
        \begin{enumerate}[leftmargin=*, topsep=0pt]
          \item[3a.] MAC đã tồn tại: trả về cấu hình hiện có mà không
            tạo bản ghi mới
          \item[1a.] Wi-Fi không kết nối được: ESP32 retry sau 5 giây
        \end{enumerate}
      \end{minipage} \\
    \textbf{Hậu điều kiện} &
      \begin{minipage}[t]{\linewidth}
        \begin{enumerate}[leftmargin=*, topsep=0pt]
          \item Node được đăng ký trong hệ thống với \texttt{api\_key}
            duy nhất
          \item Node sẵn sàng gửi dữ liệu lên topic
            \texttt{device/\{api\_key\}/data}
        \end{enumerate}
      \end{minipage} \\
    \bottomrule
  \end{tabularx}
\end{table}

\subsection{Đặc tả use case UC02 - Thu thập và lưu trữ dữ liệu
cảm biến}
\label{subsec:survey:uc02}

\begin{table}[H]
  \centering
  \caption{Đặc tả use case UC02: Thu thập và lưu trữ dữ liệu cảm biến}
  \label{tab:uc02}
  \begin{tabularx}{\textwidth}{p{4cm} X}
    \toprule
    \textbf{Mã use case}    & UC02 \\
    \midrule
    \textbf{Tên}            & Thu thập và lưu trữ dữ liệu cảm biến \\
    \textbf{Tác nhân}       & Node IoT (ESP32), Hệ thống tự động \\
    \textbf{Mục đích}       & Đọc dữ liệu môi trường định kỳ và lưu
                              trữ vào cơ sở dữ liệu \\
    \textbf{Tiền điều kiện} &
      \begin{minipage}[t]{\linewidth}
        \begin{enumerate}[leftmargin=*, topsep=0pt]
          \item Node đã đăng ký thành công (UC01)
          \item Kết nối MQTT giữa node và broker ổn định
        \end{enumerate}
      \end{minipage} \\
    \textbf{Luồng chính}    &
      \begin{minipage}[t]{\linewidth}
        \begin{enumerate}[leftmargin=*, topsep=0pt]
          \item Mỗi 5 giây, ESP32 đọc nhiệt độ, độ ẩm từ DHT11 và sinh AQI mô phỏng
          \item Đóng gói thành JSON:
            \texttt{\{"temperature": x, "humidity": y, "aqi": z\}}
          \item Publish lên topic \texttt{device/\{api\_key\}/data}
          \item Aedes broker xác thực \texttt{api\_key}, lấy
            \texttt{device\_id} từ PostgreSQL
          \item Lưu bản ghi vào bảng \texttt{readings} với timestamp
          \item Gọi \texttt{alertService} kiểm tra ngưỡng
          \item Emit \texttt{new\_reading} qua Socket.io
        \end{enumerate}
      \end{minipage} \\
    \textbf{Luồng thay thế} &
      \begin{minipage}[t]{\linewidth}
        \begin{enumerate}[leftmargin=*, topsep=0pt]
          \item[4a.] \texttt{api\_key} không hợp lệ: bỏ qua message,
            ghi log cảnh báo
          \item[1a.] DHT11 lỗi đọc: bỏ qua chu kỳ này, thử lại sau
            5 giây
        \end{enumerate}
      \end{minipage} \\
    \textbf{Hậu điều kiện} &
      \begin{minipage}[t]{\linewidth}
        \begin{enumerate}[leftmargin=*, topsep=0pt]
          \item Dữ liệu được lưu bền vững trong PostgreSQL
          \item Các client kết nối nhận được dữ liệu realtime qua
            Socket.io
        \end{enumerate}
      \end{minipage} \\
    \bottomrule
  \end{tabularx}
\end{table}

\subsection{Đặc tả use case UC03 - Xem dự báo chuỗi thời gian và
heatmap không gian}
\label{subsec:survey:uc03}

\begin{table}[H]
  \centering
  \caption{Đặc tả use case UC03: Xem dự báo và heatmap không gian}
  \label{tab:uc03}
  \begin{tabularx}{\textwidth}{p{4cm} X}
    \toprule
    \textbf{Mã use case}    & UC03 \\
    \midrule
    \textbf{Tên}            & Xem dự báo chuỗi thời gian và heatmap
                              không gian \\
    \textbf{Tác nhân}       & Người dùng (User) \\
    \textbf{Mục đích}       & Cung cấp thông tin dự báo xu hướng môi
                              trường và phân bố không gian để hỗ trợ
                              quyết định \\
    \textbf{Tiền điều kiện} &
      \begin{minipage}[t]{\linewidth}
        \begin{enumerate}[leftmargin=*, topsep=0pt]
          \item Người dùng đã đăng nhập
          \item Có ít nhất 30 bản ghi trong PostgreSQL cho node
            cần dự báo
          \item FastAPI service đang chạy
        \end{enumerate}
      \end{minipage} \\
    \textbf{Luồng chính}    &
      \begin{minipage}[t]{\linewidth}
        \begin{enumerate}[leftmargin=*, topsep=0pt]
          \item Người dùng chọn node và nhấn ``Xem dự báo''
          \item Frontend gọi \texttt{GET /api/forecast?device\_id=\{id\}}
          \item Backend Node.js gọi FastAPI service (cổng 8000)
          \item FastAPI đọc 30 bản ghi gần nhất, chuẩn hóa qua
            \texttt{scaler.pkl}
          \item Mô hình LSTM suy luận, trả về mảng 15 giá trị dự báo
          \item Frontend hiển thị biểu đồ dự báo kết hợp với dữ liệu
            lịch sử
          \item Người dùng chọn tab ``Heatmap'': frontend gọi
            \texttt{GET /api/spatial-forecast}
          \item FastAPI tính Kriging/IDW, trả về lưới
            40$\times$50 điểm
          \item Leaflet render heatmap màu sắc trên bản đồ phòng
        \end{enumerate}
      \end{minipage} \\
    \textbf{Hậu điều kiện} &
      Người dùng xem được dự báo 15 phút tới và bản đồ nhiệt phân bố
      không gian \\
    \bottomrule
  \end{tabularx}
\end{table}

%------------------------------------------------------------
\section{Yêu cầu phi chức năng}
\label{sec:survey:nonfunc}
%------------------------------------------------------------

Hệ thống được xây dựng cần đảm bảo các yêu cầu phi chức năng sau:

\textbf{Hiệu năng:} Độ trễ từ khi node gửi dữ liệu đến khi dashboard
cập nhật không vượt quá 2 giây trong điều kiện mạng bình thường. Thời
gian phản hồi của các API thông thường dưới 500ms. Thời gian suy luận
LSTM (15 phút dự báo) dưới 1 giây. Hệ thống hỗ trợ tối thiểu 10 node
IoT hoạt động đồng thời.

\textbf{Độ tin cậy:} Dữ liệu cảm biến phải được lưu trữ đầy đủ, không
mất mát trong điều kiện mạng ổn định. Khi node IoT mất kết nối MQTT,
thiết bị tự động thử kết nối lại (reconnect) theo cơ chế backoff. Hệ
thống cảnh báo phải gửi email trong vòng 30 giây kể từ khi phát hiện
vượt ngưỡng.

\textbf{Bảo mật:} Tất cả API endpoint yêu cầu xác thực JWT. Mật khẩu
người dùng được mã hóa bằng bcrypt trước khi lưu vào cơ sở dữ liệu.
Mỗi node IoT được định danh bằng \texttt{api\_key} duy nhất, được xác
thực tại backend trước khi chấp nhận dữ liệu.

\textbf{Khả năng sử dụng:} Giao diện dashboard phải trực quan, hiển thị
dữ liệu realtime mà không cần tải lại trang. Biểu đồ và heatmap phải
phản hồi ngay lập tức khi có dữ liệu mới. Hệ thống cảnh báo hiển thị
rõ ràng thông tin node, loại cảnh báo và thời điểm xảy ra.

\textbf{Khả năng mở rộng:} Kiến trúc phải cho phép bổ sung node cảm
biến mới mà không cần thay đổi code backend. Cơ sở dữ liệu PostgreSQL
được thiết kế để hỗ trợ truy vấn dữ liệu lịch sử hiệu quả với chỉ mục
phù hợp trên cột \texttt{timestamp} và \texttt{device\_id}.

\textbf{Chi phí:} Tổng chi phí phần cứng cho một node cảm biến (ESP32 +
DHT11 + nguồn) không vượt quá 150.000 VND, đảm bảo khả năng triển khai
mật độ cao với ngân sách hạn chế.
